\begin{enumerate}
    \item Durante el período 2015–2025, la economía guatemalteca mostró una tendencia de crecimiento sostenido en la mayoría de los indicadores monetarios y financieros analizados, siendo las remesas familiares y el numerario en circulación las variables que registraron el mayor crecimiento acumulado.
    \item Los análisis realizados evidencian una asociación positiva entre el Producto Interno Bruto (PIB) y variables como el crédito bancario, los medios de pago (M2) y las remesas familiares, lo que refleja la estrecha relación existente entre la actividad económica, la liquidez y el financiamiento del sector privado.
    \item El spread bancario en moneda nacional presentó un comportamiento más estable que el spread en moneda extranjera, lo cual sugiere que el mercado financiero interno ofrece menores fluctuaciones frente a los cambios originados por factores externos.
    \item Las remesas familiares constituyen uno de los principales motores de la liquidez y del crecimiento económico del país.
    \item Sin embargo, su elevada participación también evidencia una importante dependencia de ingresos provenientes del exterior, situación que representa un riesgo ante posibles cambios en el entorno económico internacional.
    \item La evolución del PIB permitió identificar un patrón estacional relativamente estable que únicamente fue interrumpido durante la pandemia de COVID-19, período en el que la economía experimentó una fuerte contracción seguida por un proceso gradual de recuperación y normalización.
\end{enumerate}
